\section{Experimental Methodology}
\label{sec:methodology}

The evaluation asks two questions. Where does the attribute-trie design reduce query latency relative to mature alternatives? Where, and why, does it not? Our design goals were \emph{semantic equivalence before speed}, \emph{fair physical design for every opponent}, and \emph{full disclosure of exclusions}. All harnesses, raw per-trial samples, plans, index inventories, and reports are publicly available (see the data availability statement).

\subsection{Systems and Configurations}
Table~\ref{tab:suites} lists the seven suites. Kiara~2.0.16 was used unmodified throughout. Opponents were Apache Lucene~10.5.2; SQLite~3.53.4 through the Xerial JDBC driver~3.53.4.0~\cite{gaffney2022sqlite}; and VoltDB Community Edition~10.1.1, the last AGPL release, from 2020~\cite{stonebraker2013voltdb}. The VoltDB results do not describe current commercial VoltDB.

\emph{Tuned} configurations received workload-specific physical design. Tuned SQLite used 38 indexes: 17 covering scalar indexes, 4 interval-bound indexes, 5 list owner, value, and interval indexes, 11 compound indexes for the tested conjunctions, and the primary key. It also ran a full \texttt{ANALYZE} and froze, per query, the fastest of 197 validated SQL alternatives (indexed joins, \texttt{IN}, \texttt{EXISTS}/\texttt{NOT EXISTS}, \texttt{EXCEPT}, \texttt{INTERSECT}). Tuned VoltDB used precompiled read-only single-partition procedures with replicated data, 37 verified indexes, and the fastest of 184 candidates, one of which timed out and was excluded. Frozen plans had to match final plans.

\emph{Generic} configurations retained per-attribute indexes but no workload-specific compound indexes or SQL search: 26 indexes and no \texttt{ANALYZE} for SQLite, and one predetermined SQL statement per query for VoltDB. ``Generic'' therefore does not mean unindexed. Kiara's schema and query formulations were held fixed in every suite. Kiara received no query-specific tuning, so the tuned suites deliberately favor the opponents.

\begin{table*}[t]
\centering
\caption{Benchmark Suites. Ratios in Section~\ref{sec:results} Are Opponent Latency Divided by Kiara Latency}
\label{tab:suites}
\setlength{\tabcolsep}{3.5pt}
\renewcommand{\arraystretch}{1.1}
\begin{tabular}{@{}lP{44mm}P{74mm}lP{22mm}@{}}
\toprule
Suite & Opponent configuration & Timed region (both engines unless noted) & Trials & Datasets \\
\midrule
S1 & Lucene, \texttt{ByteBuffersDirectory}, query cache off, numeric points and \texttt{LongRange} & In-process query construction, execution, materialization, ID extraction & 3 forks & Base five \\
S2 & VoltDB, client--server, parameterized ad hoc SQL, generic indexes & Client-visible: transport, execution, fetch, decoding & 3 forks & Base five \\
S3 & VoltDB, tuned & Server-side SQL execution and \texttt{VoltTable} creation inside a procedure; Kiara: search and eager result maps & 5 & Base five \\
S4 & VoltDB, generic & As S3 & 5 & Base five, wide \\
S5 & SQLite, tuned, private \texttt{:memory:} database & Prepared execution, all cursor steps, JDBC/JNI conversion, eager result maps & 5 & Base five \\
S6 & SQLite, generic; 60\,s per-call deadline for both engines & As S5 & 5 & Base five \\
S7 & VoltDB tuned, SQLite tuned and generic, Lucene & Per-pair boundaries as above; Lucene constant-score retrieval, document order & 5 & Wide \\
\bottomrule
\end{tabular}
\end{table*}

\subsection{Datasets and Workloads}
Table~\ref{tab:datasets} summarizes the fixtures. The retail fixture is a projection of the UCI Online Retail dataset~\cite{chen2015online}, licensed CC BY 4.0. It contains 3,788 genuine stock codes with historical prices, units sold, and the list of countries in which each product sold, aggregated from 525,462 qualifying transactions among 541,909 source rows. Catalog and schedule fixtures are generated from fixed seeds with skewed, nullable, list-valued, and interval-valued attributes. The wide fixtures vary the number of \emph{distinct} filtered attributes from 2 to 16 at target selectivities of 0.1\%, 1\%, and 10\%, giving 39 queries per fixture. Non-selector predicates accept 95\% of rows, and 90\% for \op{CONTAINS\_ALL}, so that width adds real, mostly broad predicates. Scalar-only controls stop at eight attributes, the number of applicable scalar fields. All engines receive identical input objects from one annotated record class with 27 attributes. Attributes irrelevant to a family are null.

The base workload comprises 60 scenarios in 12 categories: equality, numeric comparisons, lists, missing values, typed scalars, dates and times, date and time lists, stored ranges, combined conjunctions, result shaping, text, and dictionary operations. Together they exercise all eleven operators and all seventeen annotations.

\begin{table}[t]
\centering
\caption{Datasets}
\label{tab:datasets}
\setlength{\tabcolsep}{3pt}
\begin{tabular}{@{}lrP{50mm}@{}}
\toprule
Fixture & Records & Characteristics \\
\midrule
retail & 3,788 & Real UCI projection; market lists, fractional prices, titles \\
catalog & 10k, 100k & Status, region, tier, sale, supplier, prices, tags, sizes, titles \\
schedule & 10k, 100k & Dates, timestamps, date/time lists, numeric, prefixed, date, and time ranges \\
wide & 10k, 100k & 2--16 distinct attributes; 0.1/1/10\% selectivity \\
\bottomrule
\end{tabular}
\end{table}

\subsection{Correctness Oracle and Eligibility}
\label{sec:oracle}
Before warm-up, every timed scenario was checked against an engine-independent scan oracle for matched IDs, counts, and result shape. In S7, results were checked again after timing. Unexpected mismatches aborted the run instead of producing samples. Recognized semantic differences were labeled, retained, and excluded from speed aggregates:
\begin{itemize}
\item \textsc{Known\_Decimal\_Truncation}: Kiara's integer keys for fractional numbers (Section~\ref{sec:system}), accepted only when the exact truncation model predicted the mismatch.
\item \textsc{Text\_Result\_Difference}: keyword outputs that differ from the intended-token oracle.
\item \textsc{Unsupported}: the SQL adapters have no equivalent text, prediction, or correction operation. These modes were timed for neither engine, which is not a statement about the SQL engines' full-text features.
\item Limited queries can return different valid subsets; they were validated but excluded from ratios.
\item \textsc{Query\_Timeout} and \textsc{Build\_Failed}: recorded as outcomes, never as fabricated latencies.
\end{itemize}
Only queries that returned oracle-equivalent results in every trial of both engines entered a ratio.

\subsection{Protocol and Statistics}
Each trial ran in a fresh JVM. In S3--S7, each trial also ran in a fresh Linux container with a four-CPU quota, an 8\,GiB memory limit, and a 1--4\,GiB Java heap. VoltDB was restarted before every trial, data were rebuilt outside the timed region, engine order alternated, and scenario order was shuffled reproducibly. Warm-up comprised at least ten calls and 1\,s, followed by 15--101 measured calls within a 3\,s acquisition budget. S1 and S2 used three forks, at least three warm-up queries, and 5--61 samples within a 1\,s budget. Kiara and Lucene workers used JDK~21. The legacy VoltDB server required JDK~11. All suites ran on one shared, virtualized AMD EPYC 7763 host.

For each query we took the median across trials of the per-trial nearest-rank p50. The per-query ratio is opponent latency divided by Kiara latency, so values above one favor Kiara. A dataset-level ratio is the unweighted geometric mean of eligible per-query ratios~\cite{fleming1986how}. Its exploratory 95\% interval is a percentile bootstrap~\cite{efron1993introduction} over 5,000 whole-trial resamples with a fixed seed, which preserves within-trial correlation between queries. With five trials these intervals describe trial variability for a fixed workload. They do not establish generality. The seven suites recorded 1,042,491 raw latency samples, including tuning-candidate samples.
